\documentclass[10pt]{article}

\usepackage[utf8]{inputenc}

\usepackage[T1]{fontenc}

\usepackage{graphicx}

\usepackage[export]{adjustbox}

\graphicspath{ {./images/} }

\usepackage{amsmath}

\usepackage{amsfonts}

\usepackage{amssymb}

\usepackage[version=4]{mhchem}

\usepackage{stmaryrd}

\usepackage{bbold}



\begin{document}

\includegraphics[max width=\textwidth, center]{2023_07_09_794a8ccc5858d56eac6bg-1}

"Stat. neerlandica", 1977 , v. 31 , p. 1-24. Я. Ю. Никитин., ӘФФЕКТИВНОСТЬ СТАТИСТИ ЧЕСКОЙ ПРОЦЕДУ РЫ - понятие, используемое при сравнении статистич. процедур из данного класса с оштимальной. В математич. статистике понятие оптимальности статистич. процедуры выражается в терминах риска (функции риска) процедуры, к-рый, в свою очередь, непосредственно зависит от выбора функции потерь. Поәтому может оказаться, что одна и та же статистич. процедура является отень эффективной или даже оптимальной в каком-то одном смысле и мајо эфффективной в другом смысле.



Ә. с. п.- понятие, не совсем четко определенное; более точный сыысл оно приобретает в конкретных задачах математич. статистики, в таких, как, напр., статистических гипотез проверка и статистическое оценивание.



M. С. Никнлия.



ЭФФЕКТИВНЫЙ КРИТЕРИЙ - наиболее мощНый критерий в рассматриваемом классе статистич. критериев, имеющих один и тот же значимости уровень.



М. С. Никулин.



\begin{center}

\includegraphics[max width=\textwidth]{2023_07_09_794a8ccc5858d56eac6bg-1(2)}

\end{center}



ЮНГА ДИАГРАММА по р я д к а $m$ - Ғонаа таблица порядка $m$, в клетках к-рой каким-либо образом расположены числа $1,2, \ldots, m$, (см. табл.). ІО. Д. наз. с т а н д а р т ной, если в каждой ее строке и в каждом столбце числа идут в порядке возрастания. Числь всех Ю. Д. для данной таблицы Юнга $t$ порядка $m$ равно $m !$, а число стандартных Ю. Д. равно



\[

\frac{m !}{\Pi \lambda_{i j}}

\]



\begin{center}

\includegraphics[max width=\textwidth]{2023_07_09_794a8ccc5858d56eac6bg-1(1)}

\end{center}



где произведение берется по всем клеткам $c_{i j}$ таблицы $t$, а $\lambda_{i j}$ обозначает длину соответствующего крюка.



Ә. Б. Винбера



ЮНГА ПРИЗНАК - один из достаточных признаков сходимости Фурье ряда в тотке. Пусть функция $f(x)$ имеет период $2 \pi$, интегрируема по Јебегу на отрезке $[0,2 \pi], \varphi_{x}(t)=f(x+t)+f(x-t)-2 f(x)$ и в точке $x_{0}$ выполнены следующие условия: $\varphi_{x_{0}}(t) \longrightarrow 0$ при $t \longrightarrow+0$; Функция $\psi(t)=t \varphi_{x_{0}}(t)$ имеет конечную вариацию $\hat{V}(\delta) \equiv \operatorname{Var} \psi(t)$ на отрезке $[0, \delta], 0 \leqslant \delta \leqslant \delta_{0}$, где $\delta_{0}>0-$ нек-рое фиксированное тисло; $V(\delta)=O(\delta)$ при $\delta \longrightarrow+0$. Тогда ряд Фурье функции $f(x)$ в точке $x_{0}$ сходится к $f\left(x_{0}\right)$ (см. [2]). Ю. і. сильнее ЖКордана признака. Установлен $\mathrm{Y}$. Юнгом [1].



Jum.: [1] Y o u n g W. H., "Proc. Lond. Math. Soc.», 1916, v. 17, p. $195-236 ;$ [2] Б а p и Н. К., Тригонометрические ряды, М., 1961. Б. И. Голубов. ЮНГА СИММЕТРИЗАТОР - элемент $e_{d}$ групшового кольца группы $S_{m}$, определяемый Юнга диаграмлой $d$ порядка $m$ по следующему правилу. Пусть $R_{d}$ (coответственно $C_{d}$ ) - подгруппа группы $S_{m}$, состоящая из всех подстановок, переводящих каждое из чисел $1,2, \ldots$, $m$ в число, находящееся в той же строке (соответственно в том же столбце) диаграммы $d$. И пусть



\[

r_{d}=\sum_{g \in R_{d}}^{g}, c_{d}=\sum_{g \in C_{d}} \varepsilon(g) g,

\]



где $\varepsilon(g)= \pm 1$ - четность годстановки $g$. Тогда $e_{d}=$ $=c_{d} r_{d}$ (иногда определяіот $e_{d}=r_{d} c_{d}$ ).



Основное своӥство Ю. с. состоит в том, что он пропорционален неразложимому идемпотенту группової алгебры $\mathbb{C} S_{m}$. Коәффициент пропорциональности равен произведению длин всех кріоков диаграммы $d$.



ЮНГА ТАБЛИЦА п о р я д к а $m$-графическое представление разбиения $\lambda=\left(\lambda_{1}, \ldots, \lambda_{r}\right)$ натурального числа $m$ (где $\lambda_{i} \in \mathbb{Z}, \lambda_{1} \geqslant \lambda_{2} \geqslant \ldots \geqslant \lambda_{r}>0$, $\sum \lambda_{i}=m$ ). Ю. т. $t \lambda$ состопт из $m$ клеток, располагающихся в ее строках и столбцах таким образом, что в $i$-й строке находится

\includegraphics[max width=\textwidth, center]{2023_07_09_794a8ccc5858d56eac6bg-1(3)}



$\lambda_{i}$ клеток, п?ичем первые клетки всех строк находятся в одном (первом) столбце. Напр., разбиение $(6,5$, $4,4,1)$ числа 20 представляется Ю.т. (см. табл. слева).



Транспонированная Ю. т. $t_{\lambda}^{\prime}$ соответотвует сопряженному разбиению $\lambda^{\prime}=\left(\lambda_{1}^{\prime}, \ldots, \lambda_{s}^{\prime}\right)$, где $\lambda_{j}^{\prime}-$ число клеток в $j$-м столбце Ю. т. Так, в приведенном выше примере сопряженным разбиением будет $(5,4,4,4$, $2,1)$.



Каждая клетка Ю. т. определяет два множества клеток - так наз. к р о к и к о с о й к р ю к. Пусть $c_{i j}-$ клетка, находящаяся в $i$-й строке и $j$-м столбце данной Ю.т. Соответствующий ей крюк $h_{i j}$ есть множество, состоящее из клеток $c_{i l}$ с $l \geqslant j$ и $c_{k j}$ с $k \geqslant i$, а косой крюк есть наименьшее связное множество крайних клеток, содержащее последніою клетку $i$-й строки и последнюю клетку ј-го столбца. Напр., для изображенной слева Ю. т. крюк и косой крюк, определяемые клеткой $c_{22}$, имеют вид, показанный на табл. в центре и справа соответственно.



Д ли ш о й крюка (соответственно косого крюка) наз. число входящих в него клеток. Длина крюка $h_{i j}$ равна $\lambda_{i j}=\lambda_{i}+\lambda_{i}^{\prime}-i-j+1$. При удалении из Ю. т. косого крюка длины $p$ получается Ю. т. порядка $m-p$. В ы с о т о й крюка (соответственно косого крюка) наз. число строк, в которых располагаются его клетки.



Язык Ю. т. и Юнга диаграмм используется в теории представлений симметрических групп и в теории представлений классических групп. Он был предложен А. Юнгом (см. [1]).



Jum.: [1] J o u n g A., "Proc. Lond. Math. Soc.», 1901, ЮНГА TEOPEMА: каждое множество диаметра $d$ евклидова пространства $E_{n}$ принадлежкит шару из $E_{n}$ радиуса $r=d \sqrt{\frac{n}{2(n+1)}}$. Имеются аналоги и обобщения IO. т. (напр., замена евклидова расстояния на цругие метрики) (см. [2]).



Теорема доказана $\Gamma$. Юнгом [1].



Jum.: [1] J u n g H. W. E., "J. reine und angew. Math.", 1901, Bd 123 , S. $241 \frac{5}{7} 57$; [2] д а н ц e p Jl., $\Gamma$ p ю н 6 a у м Б., 1968: [3]' Ха л в и г е р Г., Д е б $\mathrm{p}$ у и н с $]$ Г., Комбинаторная геометрия плоскости, пер. с нем., М., 1965. "П. С. Солтан.





\end{document}